\subsection{注意力层的算术强度分析}

使用 FlashAttention 时，注意力层的计量：

\begin{itemize}
    \item \textbf{FLOPs}：$4BSTD$
    \item \textbf{传输字节}：$4BSD + 4BTD$
\end{itemize}

$$\text{注意力算术强度} = \frac{ST}{S + T}$$

代入两个阶段：

\begin{itemize}
    \item \textbf{Prefill}（$T = S$）：强度为 $S/2$，随序列长度增长
    \item \textbf{Generation}（$T = 1$）：强度为 $\frac{S}{S+1} < 1$——\textbf{始终内存受限}！
\end{itemize}

\begin{warningbox}{注意力层与 MLP 层的关键区别}
\begin{itemize}
    \item \textbf{MLP 层}：所有序列共享相同的权重矩阵（$W_\text{up}$, $W_\text{gate}$, $W_\text{down}$ 不依赖于 $B$），因此增大批大小可以分摊权重读取成本
    \item \textbf{注意力层}：每个序列有自己的 KV Cache（Q、K、V 都依赖于 $B$），因此\textbf{增大批大小无法提高算术强度}
\end{itemize}
这意味着注意力层在生成阶段\textbf{不可避免地}是内存受限的。
\end{warningbox}

